\documentclass{report}

\def\filename{Tut2Q10a}
\def\course{UECM2353}
\def\chapsec{Tutorial 2}
\def\date{\today} % to be changed

\def\headL{\textbf{\course :} \chapsec}
\def\headC{}
\def\headR{\date}

\def\footL{}
\def\footC{\thepage}
\def\footR{Carlos Wong 2303326}

\usepackage{amsmath, amssymb, amsfonts}

\usepackage[a4paper, hmargin=0.85in, vmargin=0.9in, includeheadfoot]{geometry}

\usepackage{fancyhdr}
\setlength{\headheight}{14pt}
\pagestyle{fancy}
\fancyhf{}

\lhead{{\large \headL}}
\chead{{\large \headC}}
\rhead{{\large \headR}}

\lfoot{\footL}
\cfoot{\footC}
\rfoot{\footR}

\usepackage{etoolbox}



\begin{document}
\paragraph{Q10} Determine whether the linear transformation $T$ is one-to-one, onto or neither.

(a) $T : \mathcal{P}_2 \to \mathbb{R}^3 $ defined by $T(a+bx+cx^2) = \begin{bmatrix}{2a-b\\a+b-3c\\c-a}\end{bmatrix}$

\bigskip \noindent
\textit{Solution.} \smallskip

Consider $T(a+bx+c^2) = \begin{bmatrix}{\alpha \\ \beta \\ \gamma}\end{bmatrix}, \; \alpha, \beta, \gamma \in \mathbb{R}$.

Then

\[\begin{bmatrix}{2a-b\\a+b-3c\\c-a}\end{bmatrix} = \begin{bmatrix}{\alpha \\ \beta \\ \gamma}\end{bmatrix} \implies 
\begin{bmatrix}{2 & -1 & 0 \\ 1 & 1 & -3 \\ -1 & 0 & 1}\end{bmatrix}\begin{bmatrix}{a \\ b \\ c}\end{bmatrix}=\begin{bmatrix}{\alpha \\ \beta \\ \gamma}\end{bmatrix}\]

\[ \left[\begin{array}{ccc|c}
2 & -1 & 0 & \alpha \\ 
1 & 1 & -3 & \beta \\ 
-1 & 0 & 1 & \gamma
\end{array}\right]
\xrightarrow{R_2 \to R_2 - \frac12 R_1}
\left[\begin{array}{ccc|c}
2 & -1 & 0 & \alpha \\ 
0 & \frac32 & -3 & \beta - \frac12 \alpha \\ 
-1 & 0 & 1 & \gamma
\end{array}\right]
\xrightarrow{R_3 \to R_3 + \frac12 R_1}
\left[\begin{array}{ccc|c}
2 & -1 & 0 & \alpha \\ 
0 & \frac32 & -3 & \beta - \frac12 \alpha \\ 
0 & -\frac12 & 1 & \gamma + \frac12 \alpha
\end{array}\right]
\]

\[ \xrightarrow{R_3 \to R_3 + \frac13 R_2}
\left[\begin{array}{ccc|c}
2 & -1 & 0 & \alpha \\ 
0 & \frac32 & -3 & \beta - \frac12 \alpha \\ 
0 & 0 & 0 & \gamma + \frac12 \alpha + \frac13 \beta - \frac16 \alpha
\end{array}\right]
=
\left[\begin{array}{ccc|c}
2 & -1 & 0 & \alpha \\ 
0 & \frac32 & -3 & \beta - \frac12 \alpha \\ 
0 & 0 & 0 & \gamma + \frac13 \alpha + \frac13 \beta
\end{array}\right]
\]

Thus, if $\begin{bmatrix}{\alpha \\ \beta \\ \gamma}\end{bmatrix} \in T[\mathcal{P}_2]$, then $\gamma + \frac13 \alpha + \frac13 \beta = 0$.
$\implies \gamma = -\frac13 \alpha - \frac13 \beta$.

Hence, $T[\mathcal{P}_2] = \left\{\begin{bmatrix}{\alpha \\ \beta \\ -\frac13 \alpha - \frac13 \beta}\end{bmatrix} \middle\vert \;  \alpha, \beta \in \mathbb{R}\right\} = \left\{\alpha\begin{bmatrix}{1 \\ 0 \\ -\frac13}\end{bmatrix} + \beta\begin{bmatrix}{0 \\ 1 \\ -\frac13}\end{bmatrix} \middle\vert \;  \alpha, \beta \in \mathbb{R}\right\} \neq \mathbb{R}^3 $ 

since dim($T[\mathcal{P}_2]$) $=2$ while dim($\mathbb{R}^3$) $=3$.

Therefore, $T$ is not onto.

Next, consider $T(a+bx+cx^2) = \begin{bmatrix}{0 \\ 0 \\ 0}\end{bmatrix}$.

From the rank-nullity theorem, we can see that dim$(ker(T))$ $= 3 - \text{dim}(T[\mathcal{P}_2]) = 3 - 2 = 1$.

Hence, there are at least two distinct quadratic polynomials $p(x)$ and $q(x)$ s.t. $T(p(x)) = T(q(x)) = \begin{bmatrix}{0 \\ 0 \\ 0}\end{bmatrix}$.

Since $p(x) \neq q(x)$, $T$ is not one-to-one.

Hence, $T$ is neither one-to-one nor onto. \hfill $\square$

\end{document}
